\documentclass[11pt,reqno]{amsart}
\usepackage[margin=1.12in]{geometry}
\usepackage{amssymb}
\usepackage{mathtools}
\usepackage{microtype}
\usepackage[hidelinks]{hyperref}
\hypersetup{pdftitle={Asymptotic counts for powerful triples with arbitrary common divisor},pdfauthor={Roy van Rijn}}

\newtheorem{theorem}{Theorem}[section]
\newtheorem{proposition}[theorem]{Proposition}
\newtheorem{lemma}[theorem]{Lemma}
\newtheorem{corollary}[theorem]{Corollary}
\newtheorem{hypothesis}[theorem]{Hypothesis}
\DeclareMathOperator{\rad}{rad}
\newcommand{\eps}{\varepsilon}
\newcommand{\sord}{\boldsymbol{s}}

\title[Asymptotic counts for powerful triples]{Asymptotic counts for powerful triples with arbitrary common divisor}
\author{Roy van Rijn}
\urladdr{\url{https://orcid.org/0009-0000-2808-0645}}
\subjclass[2020]{11D45, 11N37}
\keywords{powerful integers, additive triples, asymptotic expansions, primitive quotients}

\begin{document}
\begin{abstract}
We count ordered positive triples \(a+b=c\le X\) with prescribed fullness
orders and arbitrary common divisor. Assuming a stated bound for primitive
quotients with a small product of deficient primes, we obtain an expansion in powers
\(X^{1/(M+j)}\), where \(M\) is the largest prescribed order, for sufficiently
large comparable fixed orders. The coefficients have strictly alternating
signs, and the positive leading coefficient is an absolutely convergent sum over
primitive quotients. The expansion holds in \(\ell^1\) over those quotients
and gives a limiting law with a power-saving error in total variation.
The contribution from quotients with a large product of deficient primes
has an unconditional expansion. The proof combines a coprimality-uniform
full-integer count with the cost of restoring divisibility after taking out
the common divisor.
\end{abstract}
\maketitle

\section{Introduction and main results}

For an integer \(s\ge2\), a positive integer is \(s\)-full if every prime
that divides it occurs to exponent at least \(s\). In particular, \(1\) is
\(s\)-full. Fix three orders \(\sord=(s_1,s_2,s_3)\), and put
\[
  m=\min_i s_i,\qquad M=\max_i s_i.
\]
Let \(A_{\sord}(X)\) count ordered positive triples \(a+b=c\le X\) whose
entries are respectively \(s_1\)-, \(s_2\)-, and \(s_3\)-full. Dividing by
\(g=\gcd(a,b,c)\) gives a unique primitive quotient \(q=(u,v,w)\), with
\(u+v=w\) and \(\gcd(u,v,w)=1\). Let \(a_q(X)\) count the triples with
quotient \(q\).

The quotient \(1+1=2\) explains the scale of the count. For any
\(M\)-full odd integer \(b\), the common divisor \(2^M b\) gives a triple
of the required orders with largest entry \(2^{M+1}b\). Taking \(b\) to
range over odd \(M\)-th powers already gives \(\gg_M X^{1/M}\) triples.
In general, the part of \(g\) supported on primes of \(uvw\) must restore
fullness at certain primes. The remaining part of \(g\) varies over
\(M\)-full integers coprime to \(uvw\). This second part produces the
asymptotic terms. The cost of the first part can be used to show that
the resulting series over all primitive quotients converges.

Browning and Verzobio \cite[Theorem~1.1]{BV} count primitive triples of
fixed nearby orders. For \(s_i=k\) their upper-bound exponent is
\(2/k-1/k^2+O(k^{-5/2})\), apart from the arbitrary exponent allowance.
Their theorem fixes the common divisor to one. A primitive quotient in the
present count can have a prime of any positive valuation in its entries.
On the one-integer side, Cohen and Davis \cite[\S\S6--7]{CD} give two-term
estimates for the number of \(M\)-full integers. We use a related expansion
with a variable coprimality restriction and an error that can be summed
over primitive quotients.

Write \((q_1,q_2,q_3)=(u,v,w)\) and call a prime \(p\) deficient for
\(q\) if \(0<v_p(q_i)<s_i\) for some \(i\). Let \(D(q)\) be the product
of the deficient primes, with an empty product equal to one. For powers
of two \(Y,H\), let \(Q(Y,H)\) count primitive quotients satisfying
\(Y\le w<2Y\) and \(H\le D(q)<2H\).

\begin{hypothesis}[Primitive quotient count]\label{hyp:quotient}
For the prescribed orders there is \(0<\kappa<1\) such that, for every
\(\eps>0\),
\begin{equation}\label{eq:quotient}
  Q(Y,H)\ll_{\sord,\eps}Y^{(1-\kappa)/M+\eps}H^{m/M}
   \qquad (Y\ge2,\ 1\le H\le Y^{20/M}).
\end{equation}
\end{hypothesis}

The hypothesis is needed only for the displayed range of \(H\). A simple
radical count suffices above it. Section~\ref{sec:endpoint} derives
\eqref{eq:quotient} from two further hypotheses about small-budget
\(abc\) triples.

\begin{theorem}\label{thm:main}
Assume Hypothesis~\ref{hyp:quotient}, \(M/m\le 11/10\), and let \(J\ge0\) be
an integer with
\begin{equation}\label{eq:range}
  M\ge 16(J+1)^2,\qquad M\ge 100(J+1)/\kappa.
\end{equation}
Then there are absolutely summable real functions \(\mu_j\) on the ordered
primitive quotients, with \((-1)^j\mu_j(q)>0\) for every \(q\), such that,
for every \(\eps>0\),
\begin{equation}\label{eq:main}
  \sum_q\left|a_q(X)-\sum_{j=0}^{J}\mu_j(q) X^{1/(M+j)}\right|
   \ll_{\sord,J,\eps} X^{1/(M+J+1)+\eps}\qquad (X\ge 1).
\end{equation}
In particular
\begin{equation}\label{eq:scalar}
\begin{aligned}
  A_{\sord}(X)&=\sum_{j=0}^J C_{\sord,j}X^{1/(M+j)}
       +O_{\sord,J,\eps}\bigl(X^{1/(M+J+1)+\eps}\bigr),\\
  C_{\sord,j}&=\sum_q\mu_j(q),\quad (-1)^jC_{\sord,j}>0.
\end{aligned}
\end{equation}
\end{theorem}

The leading exponent is \(1/M\), including when the orders differ. For
equal orders \(k\), this gives \(A_{(k,k,k)}(X)\sim C_{(k,k,k),0}X^{1/k}\).
For a fixed \(\kappa>0\) the admissible number of terms grows at least
as a constant multiple of \(\sqrt M\) as \(M\) tends to infinity. The
orders are fixed when taking the limit in \(X\).

\begin{corollary}\label{cor:law}
Assume the hypotheses of Theorem~\ref{thm:main} with \(J=0\). For sufficiently
large \(X\), choose a triple counted by \(A_{\sord}(X)\) uniformly and let
\(\mathbb P_X\) be the law of its primitive quotient. There is a limiting law
\(\mathbb P_\infty(q)=\mu_0(q)/C_{\sord,0}\), and
\[
  \|\mathbb P_X-\mathbb P_\infty\|_{\mathrm{TV}}
  \ll_{\sord,\eps} X^{-1/(M(M+1))+\eps}.
\]
There is \(\delta>0\), depending on the orders, such that
\(\mathbb P_X(w>T)+\mathbb P_\infty(w>T)\ll_{\sord}T^{-\delta}\)
uniformly for sufficiently large \(X\) and all \(T\ge1\).
\end{corollary}

\begin{proposition}\label{prop:sector}
Suppose \(M/m\le11/10\), \(M>20\), and \(M\ge16(J+1)^2\). The expansion
\eqref{eq:main} holds unconditionally with the sum restricted to
\(q\) with \(D(q)>w^{20/M}\). The corresponding scalar expansion has
strictly alternating coefficients and a positive leading coefficient.
\end{proposition}

\section{The common divisor}\label{sec:divisor}

Put \(P(q)=\rad(uvw)\). The three entries of a primitive additive quotient
are pairwise coprime, so a prime of \(P(q)\) divides exactly one entry.
Split the common divisor uniquely as \(g=fb\), where all primes of \(f\)
divide \(P(q)\) and \((b,P(q))=1\). For \(p\mid P(q)\) put
\[
  K_p=\max_i(s_i-v_p(q_i)).
\]
At least two valuations \(v_p(q_i)\) are zero, so \(K_p\ge m\). A positive
valuation of \(f\) at \(p\) can be any integer at least \(K_p\). Valuation
zero is allowed precisely when \(p\nmid D(q)\). Let \(\mathcal F(q)\)
be the resulting set of positive integers \(f\). At a prime outside
\(P(q)\), all three quotient valuations are zero. Hence the coprime part
\(b\) can be any \(M\)-full integer coprime to \(P(q)\). Writing
\(B_M(U;P)\) for the number of those integers at most \(U\), we have
\begin{equation}\label{eq:fibres}
  a_q(X)=\sum_{f\in\mathcal F(q)} B_M\bigl(X/(wf);P(q)\bigr).
\end{equation}
For the quotient \(1+1=2\), for example, \(K_2=\max(s_1,s_2,s_3-1)\),
and \(\mathcal F(q)=\{2^e:e\ge K_2\}\). This gives a full description of
the triples in that quotient, including the family in the introduction.

For every \(\sigma>0\) and \(\eps>0\) the Euler product of the local sums
gives
\begin{equation}\label{eq:repair}
  \sum_{f\in\mathcal F(q)} f^{-\sigma}
    \le D(q)^{-m\sigma}\prod_{p\mid P(q)}(1-p^{-\sigma})^{-1}
    \ll_{\sigma,\eps} D(q)^{-m\sigma}w^\eps.
\end{equation}
For the last inequality, all but finitely many of the local factors are
at most \(p^{\eps/3}\), and \(P(q)\le w^3\). The same comparison gives
\(C^{\omega(P(q))}\ll_{C,\eps} w^\eps\) for any fixed \(C\ge1\), where
\(\omega(P)\) is the number of distinct prime divisors of \(P\).

Set \(\lambda_0=\sum_i s_i^{-1}\). At a nondeficient prime the valuation
already pays for its contribution to the radical. Thus
\begin{equation}\label{eq:radical}
  P(q)\le D(q)\prod_i q_i^{1/s_i}\le D(q)w^{\lambda_0}.
\end{equation}
This gives the unconditional bound
\begin{equation}\label{eq:two}
 Q(Y,H)\ll_{\sord,\eps}Y^{2\lambda_0/3+\eps}H^{2/3}.
\end{equation}
Indeed, the two entries with smallest radicals have radical product at
most \(P(q)^{2/3}\). Those two radicals are coprime, so fixing their
product leaves at most \(2^{\omega(P')}\) assignments, where \(P'\) is
the product. There are three choices of the omitted entry and the two
retained values determine it. For a fixed radical \(R\) and \(n\le 2Y\),
Rankin's inequality bounds the number of possible values of \(n\) by
\((2Y)^\alpha\prod_{p\mid R}(1-p^{-\alpha})^{-1}\) for every
\(\alpha>0\). If \(R\le 2Y\), the product is \(O_{\alpha,\beta}(Y^\beta)\)
for each \(\beta>0\) by the same local comparison as in
\eqref{eq:repair}. Summing over \(P'\ll_{\sord}Y^{2\lambda_0/3}H^{2/3}\)
then proves \eqref{eq:two}, since \(H\ll Y^3\) and all allowances can
be chosen arbitrarily small. This is the two-smallest-radicals argument
used in \cite[Proposition~1.1]{BBLT}, here applied to \eqref{eq:radical}.

\section{Full integers subject to a coprimality restriction}\label{sec:single}

For squarefree \(P\) define
\[
  \rho_P=\prod_{p\mid P}(1-p^{-1}),\qquad
  \zeta_P(t)=\zeta(t)\prod_{p\mid P}(1-p^{-t}).
\]
Set \(a_j=M+j\) for \(0\le j\le J\) and, initially for real \(t>1/M\),
define \(G(t;P)\) by
\begin{equation}\label{eq:euler}
  \prod_{p\nmid P}\left(1+\sum_{\ell\ge M} p^{-\ell t}\right)
   =G(t;P)\prod_{j=0}^J\zeta_P(a_jt).
\end{equation}

\begin{lemma}\label{lem:single}
If \(M\ge16(J+1)^2\) and \(1/(M+J+1)<\sigma<1/(M+J)\), then uniformly
for squarefree \(P\) and \(U>0\),
\begin{align}
 B_M(U;P)&=\sum_{j=0}^J b_j(P)U^{1/a_j}
        +O_{M,J,\sigma}\bigl(2^{(J+1)\omega(P)}U^\sigma\bigr),\label{eq:one}\\
 b_j(P)&=\rho_P G(1/a_j;P)\prod_{i\ne j}\zeta_P(a_i/a_j).
 \label{eq:b}
\end{align}
The coefficients are bounded uniformly in \(P\), satisfy
\((-1)^j b_j(P)>0\), and the coefficient with index \(j\) is independent
of the truncation index \(J\ge j\).
\end{lemma}

\begin{proof}
At a prime outside \(P\), multiplying the factor on the left of
\eqref{eq:euler} by \(\prod_{j=0}^J(1-p^{-a_jt})\) cancels every
positive exponent below \(M+J+1\), since \(M>J\). The product defining
\(G(t;P)\) therefore has absolutely summable Dirichlet coefficients at
each real \(\sigma>1/(M+J+1)\), uniformly in \(P\). To see this at the
coefficient level, the sum of absolute values of the nonconstant local
coefficients weighted by \(p^{-\ell\sigma}\) is
\(O_{M,J,\sigma}(p^{-(M+J+1)\sigma})\), from the geometric series in
\eqref{eq:euler}. At each real point \(1/a_j\) all factors in the Euler
product for \(G\) are positive, and absolute convergence gives a strictly
positive product.

We first count product powers. For a nonempty subset \(I\subseteq
\{0,\ldots,J\}\), write \(d=|I|\), \(A_I=\sum_{i\in I}a_i\), and
\(e_I=(d-1)(1+J/M)/A_I\). For \(U\ge1\), the number \(F_I(U;P)\) of
tuples of positive integers \(x_i\) coprime to \(P\) with
\(\prod_{i\in I}x_i^{a_i}\le U\) satisfies
\begin{equation}\label{eq:hyperbola}
 F_I(U;P)=\rho_P\sum_{i\in I}U^{1/a_i}\prod_{j\in I\setminus\{i\}}
           \zeta_P(a_j/a_i)
      +O_{M,J}\bigl(2^{d\omega(P)}U^{e_I}\bigr).
\end{equation}
For \(d=1\), this follows by inclusion--exclusion in the coprime
count. For \(d\ge2\), let \(V=U^{1/A_I}\). Some coordinate is at most
\(V\), so use inclusion--exclusion on those coordinates. For a nonempty
proper subset \(K\) of coordinates constrained by \(x_i\le V\), apply
induction to \(I\setminus K=B\) at height \(U/\prod_{j\in K}x_j^{a_j}\).
If \(b=|B|\), then for every \(j\in K\),
\[
  a_j e_B\le (1-1/b)(1+J/M)^2<1.
\]
Indeed \((1+J/M)^2<1+1/(J+1)\) under the size condition for \(M\).
Summing the inductive errors over the small coordinates therefore gives
\[
 O_{M,J}\bigl(2^{b\omega(P)}V^{|K|+(b-1)(1+J/M)}\bigr),
\]
which fits the error in \eqref{eq:hyperbola}.

For the main terms use, for fixed real \(t>0\), \(t\ne1\),
\begin{equation}\label{eq:weighted}
  \sum_{\substack{x\le V\\(x,P)=1}}x^{-t}
  =\zeta_P(t)+\rho_P\frac{V^{1-t}}{1-t}
      +O_t(2^{\omega(P)}V^{-t}).
\end{equation}
Partial summation proves the unrestricted formula, with the zeta
constant defined by continuation for \(0<t<1\). Apply it to
\(V/d'\) for divisors \(d'\mid P\) with \(d'\le V\), multiplying by
\(\mu(d')d'^{-t}\), where \(\mu\) is the M\"obius function. For each
omitted \(d'>V\), both corresponding main terms are
\(O_t(V^{-t})\). This proves the stated uniform error.

Fix a pole index \(i\), and for \(j\ne i\) in \(I\) write
\(Z_j=\zeta_P(a_j/a_i)\) and
\(R_j=\rho_P V^{1-a_j/a_i}/(1-a_j/a_i)\). The alternating sum over
nonempty \(K\subseteq I\setminus\{i\}\) of products with \(Z_j+R_j\)
on \(K\) and \(Z_j\) elsewhere equals
\(\prod_{j\ne i} Z_j+(-1)^d\prod_{j\ne i}R_j\). The first product gives
the corresponding main term in \eqref{eq:hyperbola}. The second
product, after multiplication by \(\rho_P U^{1/a_i}\), contributes
\[
  (-1)^d\rho_P^d V^d\frac{a_i^{d-1}}{\prod_{j\ne i}(a_i-a_j)}.
\]
The fractions sum to one by comparing leading coefficients in the
Lagrange interpolation formula for \(z^{d-1}\). The resulting term
cancels the leading contribution \((-1)^{d+1}\rho_P^dV^d\) from the
box with all coordinates at most \(V\). That box has error
\(O_{M,J}(2^{d\omega(P)}V^{d-1})\), by the coprime single-variable
count in each slot. In a product using an error from \eqref{eq:weighted},
one slot contributes \(V^{-a_j/a_i}\), canceling its share of the initial
\(V^{A_I/a_i}\). The initial share from \(i\) is \(V\). Each of the
other slots contributes at most \(V^{1+J/M}\) after combining it with
its initial share, since either \(Z_j\) is bounded or \(R_j\) brings
\(V^{1-a_j/a_i}\). This gives the error of \eqref{eq:hyperbola}.

For \(I=\{0,\ldots,J\}\) we have \(e_I<1/(M+J+1)\), by
\((1-1/(J+1))(1+(J+1)/M)^2<1\). Convolve \eqref{eq:hyperbola} with
the coefficients of \(G\). Extending each main coefficient to include
the terms with the convolution variable greater than \(U\) costs
\(O_{M,J,\sigma}(U^\sigma)\) by absolute convergence at \(\sigma\).
The summed errors have the bound in \eqref{eq:one} since \(e_I<\sigma\).
Uniform boundedness of \(b_j(P)\) follows from absolute convergence of
\(G\) and \(|\zeta_P(t)|\le|\zeta(t)|\) for fixed real \(t>0\),
\(t\ne1\). For \(U<1\) the count is zero, and each proposed term is
\(O(U^\sigma)\). Exactly \(j\) factors in the product in \eqref{eq:b}
have arguments in \((0,1)\), where \(\zeta\) is negative by its
alternating series representation. All other factors are positive.
Finally, increasing \(J\) changes \(G(t;P)\) by dividing by the added
\(\zeta_P(a_i t)\). At an existing evaluation point \(1/a_j\), the added
factors are finite and cancel in \eqref{eq:b}.
\end{proof}

\section{Summation over primitive quotients}\label{sec:sum}

Choose \(\sigma\) just above \(1/(M+J+1)\) and below
\(1/(M+J)\). We first show that for some \(\delta>0\),
\begin{equation}\label{eq:majorant}
 \sum_{q:\,w>T} 2^{(J+1)\omega(P(q))}w^{-\sigma}
                 \sum_{f\in\mathcal F(q)} f^{-\sigma}
    \ll_{\sord,J,\sigma}T^{-\delta}\qquad(T\ge1).
\end{equation}
For \(w\) and \(D(q)\) in their respective dyadic bins \(Y,H\), the
contribution is bounded by the bin count times
\(O_{\sord,J,\sigma,\eps}(Y^{-\sigma+\eps}H^{-m\sigma})\) by
\eqref{eq:repair}. Write \(h=\log H/\log Y\). For \(h\le20/M\),
Hypothesis~\ref{hyp:quotient} bounds the exponent of \(Y\) by
\[
 -\frac{\kappa}{M}+(1/M-\sigma)(1+mh)+O(\eps)
 \le -\frac{\kappa}{M}+\frac{21(J+1)}{M^2}+O(\eps).
\]
This is negative with a fixed margin at fixed orders by \eqref{eq:range}.
For \(h>20/M\), use \eqref{eq:two}. The exponent becomes at most
\[
  \frac{2\lambda_0}{3}-\sigma-(m\sigma-2/3)h+O(\eps).
\]
Here \(\lambda_0\le 33/(10M)\) and
\(m\sigma>(10/11)(16/17)>5/6\), by the order ratio and the bound
\(M\ge16(J+1)^2\). These estimates give a fixed negative margin
of order \(1/M\), since \(2\lambda_0/3\le 11/(5M)\) whereas
\((m\sigma-2/3)h>10/(3M)\). The number of \(H\)-bins is
\(O(1+\log Y)\) because \(D(q)\le w^3\). Taking the exponent
allowances sufficiently small and summing the resulting geometric
bound over dyadic \(Y\) proves \eqref{eq:majorant}.

By \eqref{eq:fibres} and Lemma~\ref{lem:single} the coefficient at an
individual quotient is
\begin{equation}\label{eq:mu}
 \mu_j(q)= b_j(P(q)),w^{-1/a_j}
                    \sum_{f\in\mathcal F(q)}f^{-1/a_j}.
\end{equation}
The local exponent ranges defining \(\mathcal F(q)\) are nonempty. The
uniform bound for \(b_j(P)\), together with \(\sigma<1/a_j\) and
\eqref{eq:majorant}, gives absolute convergence and the tail
\(\sum_{w>T}|\mu_j(q)|\ll T^{-\delta}\). The sign of each weight is
\((-1)^j\) by Lemma~\ref{lem:single}. The triangle inequality, applied
first to each \(f\)-summation, bounds the left side of \eqref{eq:main}
by \(X^\sigma\) times the convergent series in \eqref{eq:majorant}
at \(T=1\). We can choose \(\sigma<1/(M+J+1)+\eps\), completing the
proof of Theorem~\ref{thm:main}.

For Proposition~\ref{prop:sector}, restrict to \(D(q)>w^{20/M}\).
The same second exponent estimate applies, with
\(h\ge20/M-\log 2/\log Y\) in the dyadic bins. The difference costs
only a constant factor in the summation. The quotient \(1+1=2\) lies
in this sector when \(M>20\). This proves nonvanishing of all the
sector coefficients from \eqref{eq:mu}.

For Corollary~\ref{cor:law}, take \(J=0\) in \eqref{eq:main}, and divide
by \(A_{\sord}(X)\asymp_{\sord}X^{1/M}\). The error made by replacing
\(X^{1/M}/A_{\sord}(X)\) with \(1/C_{\sord,0}\) is bounded in
\(\ell^1\) by the same order, since \(\sum_q\mu_0(q)=C_{\sord,0}\).
This gives the total-variation estimate. Restricting the sum of
remainders to \(w>T\), its bound becomes \(O(X^\sigma T^{-\delta})\)
by \eqref{eq:majorant}. Use that bound and the coefficient tail from
\eqref{eq:mu}, then divide by \(A_{\sord}(X)\) to obtain the asserted
uniform tail.

\section{The small-budget counting input}\label{sec:endpoint}

The upper bound \eqref{eq:quotient} can be reduced to two statements about
primitive additive triples of small radical. To state the second, use
boxes of factorizations \(q_i=c_i\prod_{e=1}^E x_{i,e}^e\) with positive
\(c_i\) fixed within each box, pairwise coprime roots within each entry, and primitive
totals satisfying \(q_1+q_2=q_3\) and \(Y\le q_3<2Y\). The integer
\(x_{i,e}\) lies in \([B_{i,e},2B_{i,e})\), where \(B_{i,e}\) is a power
of two. Define the count exponent \(d_0\) as the logarithm of the number of
solutions divided by \(\log Y\). The logarithmic root sizes are
\[
 L_{i,e}=\frac{\log B_{i,e}}{\log Y},\qquad L=\sum_{i,e}L_{i,e}.
\]
We need the following inputs.

\begin{hypothesis}[Small-budget counts and endpoint profiles]\label{hyp:endpoint}
\leavevmode
\begin{enumerate}
\item For every \(\eta>0\) and all sufficiently small fixed \(\lambda>0\),
the number of ordered positive primitive \(u+v=w\le X\) with
\(\rad(uvw)<w^\lambda\) is
\(O_{\lambda,\eta,\eps}(X^{(1/3+\eta)\lambda+\eps})\) for every \(\eps>0\).
\item A profile is obtained by fixing \(E\) and taking height limits of
\(d_0\), \(L_{i,e}\), and \(\log c_i/\log Y\) through populated boxes
as \(Y\to\infty\). The cutoff \(E\) may vary between profiles. Retain
the same notation for the limiting exponents and the total root mass \(L\). Along
any sequence of the resulting profiles with \(L>0\) tending to zero and
total coefficient exponent \(o(L)\), one has \(\limsup d_0/L\le1/3\).
When \(d_0/L\to1/3\), every simultaneous
weak limit of the finite measures
\begin{equation}\label{eq:measures}
  \nu_i=\sum_e\frac{L_{i,e}}{L}\,\delta_{eL}
\end{equation}
has total nonatomic mass at least \(1/2\) summed over the three entries.
Here \(\delta_r\) is unit point mass at \(r\).
\end{enumerate}
\end{hypothesis}

The exponent \(1/3\) in the second input alone would permit three point
masses at scaled degree three. The nonatomic mass condition rules out
that profile, which is forced when the quotients are almost full.

\begin{proposition}\label{prop:endpoint}
Hypothesis~\ref{hyp:endpoint} implies that there are absolute
\(\kappa,\theta>0\) such that \eqref{eq:quotient} holds for all sufficiently
large fixed orders with \(M/m\le1+\theta\). In fact the same bound holds
for all \(H\ge1\).
\end{proposition}

\begin{proof}
Suppose no such constants exist. For each integer \(n\ge2\) one can
choose a triple of fixed orders with \(m\ge n\), \(M/m\le1+1/n\), for
which the assertion with \(\kappa=1/n\) fails for some \(\eps>0\).
For each fixed triple of orders there are unbounded heights in dyadic
bins witnessing the failure. Select a height subsequence on which
\(\log H/\log Y\to h\). This is possible by \(D(q)\le w^3\).
The lower count exponent along that subsequence is at least
\((1-1/n)/M+(m/M)h\).

By \eqref{eq:two},
\((m/M-2/3)h\le (2/3)\lambda_0-(1-1/n)/M\), so \(h=O(1/M)\) as
\(n\to\infty\). By \eqref{eq:radical}, any fixed budget strictly
larger than \(\lambda_0+h\) applies at sufficiently large heights in
that subsequence. For a fixed \(\eta\in(0,1/3)\) and sufficiently large
\(n\), apply Hypothesis~\ref{hyp:endpoint}(1) at these fixed budgets
and decrease the budget to \(\lambda_0+h\), after taking the height
limit and letting the exponent allowance tend to zero. This gives
\[
 (m/M-1/3-\eta)Mh\le(1/3+\eta)M\lambda_0-(1-1/n).
\]
As \(n\to\infty\), \(m/M\to1\) and \(M\lambda_0\to3\). Letting
\(\eta\downarrow0\) shows that \(Mh\to0\).

For these bins we use the valuation factorization of
Bernert, Browning, Lichtman and Ter\"av\"ainen \cite[Lemma~2.2]{BBLT}.
It preserves prime valuations through a cutoff \(E\) in their own
coordinates. A valuation \(j>E\) is written as
\(j=k\lfloor j/k\rfloor+r\) with \(0\le r<k\), and contributes to
coordinate \(k\) and to the exceptional coefficient \(c_i\). For an
arbitrarily small fixed \(\rho>0\), their construction takes
\(k=\lceil1/\rho\rceil\) and \(E=6\lceil\rho^{-2}\rceil\), and gives
\(c_i\le q_i^\rho\). Thus \(\rho\) can be taken sufficiently small at
each fixed triple of orders that \(k\ge M\) and \(\rho=o(1/M)\) as
\(n\to\infty\). Every root at an index below its entry's required order
then consists of primes occurring to exactly that index in the quotient.
In particular the product of these roots equals \(D(q)\). This equality
includes the case in which a letter has no deficient primes.

Freeze the three coefficients and split the roots dyadically, keeping
one box with at least the average count. The coefficient choices cost
\(Y^{3\rho+o(1)}\), and the number \(O((1+\log Y)^{3E})\) of dyadic
boxes costs \(Y^{o(1)}\) at fixed \(n\). Along a further height
subsequence take limits of all the box exponents. In the resulting
profile,
\begin{align*}
 d_0&\ge (1-1/n)/M+(m/M)h-3\rho,\\
 L&\le\lambda_0+h+3\rho,\qquad \sum_i\sum_{e<s_i} L_{i,e}\le h.
\end{align*}
The total root mass bound follows from \eqref{eq:radical}, since the
product of the roots is bounded by the radical times the product of the
exceptional coefficients. This last comparison follows prime by prime
from the valuation construction, or from \cite[Lemma~2.2]{BBLT}.
Counting all root tuples gives \(d_0\le L\) in the height limit. Thus
\(L\gg1/M\), and the total coefficient exponent is at most
\(3\rho=o(L)\) as \(n\to\infty\). The exponent bound
in Hypothesis~\ref{hyp:endpoint}(2) now forces \(ML\to3\) and
\(d_0/L\to1/3\).

The measures \eqref{eq:measures} have first moments at most one
by \(q_i\le2Y\) and the height limit. They are tight,
and their total mass, summed over \(i\), is one. Their mass at indices
below a required order tends to zero by the bound for the low roots.
At every other index \(eL\ge s_i L\to3\). Every simultaneous weak
limit is therefore supported on \([3,\infty)\). The moment bound
gives each limiting measure mass at most \(1/3\), and hence each has
mass exactly \(1/3\). The same moment bound then forces all three
measures to be \(\tfrac13\delta_3\). This contradicts the nonatomic
mass condition and proves the proposition.
\end{proof}

\paragraph{AI assistance and review.}
ChatGPT (OpenAI) assisted with mathematical arguments, literature checks,
exposition, proof review and verification code.

\begin{thebibliography}{9}
\bibitem{BBLT}
C.~Bernert, T.~Browning, J.~D. Lichtman, and J.~Ter\"av\"ainen,
\emph{Bounds on the exceptional set in the \(abc\) conjecture},
\href{https://arxiv.org/abs/2410.12234v2}{arXiv:2410.12234v2}, 2026.

\bibitem{BV}
T.~Browning and M.~Verzobio, \emph{Sums of three powerful numbers},
\href{https://arxiv.org/abs/2608.24512v1}{arXiv:2608.24512v1}, 2026.

\bibitem{CD}
E.~Cohen and K.~J. Davis, \emph{Elementary estimates for certain types of
integers}, Acta Sci. Math. (Szeged) \textbf{31} (1970), 363--371.
\href{https://acta.bibl.u-szeged.hu/14284/1/math_031_fasc_003_004_363-371.pdf}{Journal archive}.
\end{thebibliography}
\end{document}
